\section{结论}

本文基于 ADRR 派生开放航空数据，对多时域区域航空轨迹预测进行了大规模实验研究。以层级 H3 GRU 基线为起点，我们引入 horizon-query Transformer decoder，并进一步加入显式多尺度时间分支。最终模型相较 GRU 基线和单尺度 Transformer 消融模型，在 top-k 准确率、NLL、Brier 分数和校准性上均取得提升，同时生成更小的候选区域集合。经典 Markov 基线明显较弱，这说明该任务需要比简单经验转移建模更强的时间表示学习能力。

H3 refinement 实验进一步说明，候选召回和校准性应当一起报告。一个激进扩展候选区域的模型，在只看召回时可能显得更强，但其概率可信度会下降。因此，大规模区域航空预测应被评估为准确性、不确定性质量和空间效率的联合问题。

总体而言，实验结果支持将层级 H3 预测与多尺度 horizon-query 时间编码结合，作为开放轨迹档案上长时域区域航空预测的强基线。它也为下一步方法研究提供了稳定基础：在缺失和不规则飞行观测下，进一步构建显式不确定性感知的区域预测模型。
